% Ad Familiares 14.5 (ad-familiares-14-05) --- English translation
% Translated by Claude (Anthropic) from the Latin (Perseus).
% Style guide: STYLE.md.

\ciceroLetterOpener{TVLLIVS S. D. TERENTIAE SVAE.}

\ciceroSection{1} If you and Tullia, light of my life, are well, then I and our dearest Cicero are well. We came to Athens on the day before the Ides of October, having had thoroughly contrary winds and a slow and unpleasant crossing. As we were stepping off the ship Acastus was on hand with letters, on the twenty-first day --- briskly enough indeed. I have received your letter, from which I gather that you have been afraid the earlier ones were never delivered to me. All of them were delivered, and you have written everything out most carefully, and that has been most welcome to me. Nor was I surprised that this letter, which Acastus brought, was a short one; for by now you are expecting me in person, or rather us, who do wish to come to you as quickly as we can, even though I see well enough what sort of public situation I am coming back to. I have learned, from the letters of many friends which Acastus has brought, that the matter is tending toward arms, so that, when I arrive, it will not be open to me to dissemble what I feel. But since one's fortune must be faced, I shall press on the more quickly to get there, so that we may the more easily deliberate on the whole business. I should like you, so far as the state of your health allows, to come as far as you possibly can to meet me.

\ciceroSection{2} On the Precius inheritance (which is in truth a great grief to me, for I loved the man dearly): I should like you to see to it, if the auction takes place before my arrival, that Pomponius --- or, if he is unable, then Camillus --- look after our business; once we have come home safe we shall handle the rest ourselves. If however you have already set out from Rome, you will see to it all the same that the matter is so handled. We hope, the gods helping, to be in Italy about the Ides of November. Take care, my sweetest and most-longed-for Terentia, if you love us, that you keep well. Farewell. Athens, the seventeenth day before the Kalends of November.
